\clearpage
\SourceLeaf{237}%
\begin{center}
\Large\textsc{[Lübsens Parallelentheorie.]}
\end{center}


\begin{center}
\emph{[\textup{\textsc{Schumacher}} an \textup{\textsc{Gauss}}. Altona, 31. December 1835.]}
\end{center}

\{\mbox{.\ .\ .\ .\ .\ .\ .\ .\ .} \textsc{Lübsen} quält mich sehr mit seiner Parallelentheorie
(wahrscheinlich tritt die Nemesis hier ein, um mir meine eigenen
Versuche zu vergelten). Ich lege Ihnen seinen Beweis bei. Ich habe ihm
vergebens bemerkt, dass das Schieben der Linien besser am
Parallellineal, als in der Theorie auszuführen sei; er bleibt doch bei
seiner Überzeugung, den geometrischen Beweis gefunden zu haben.
Vielleicht wäre es gut, da er sonst doch ein gescheuter und
bescheidener Mann ist, wenn Sie ihn mit ein Paar Worten aus dem
Irrthum rissen. Es muss wohl daran liegen, dass ich ihm seinen
Irrthum nicht so deutlich vorstellen kann, wie ich ihn selbst
erkenne. \mbox{.\ .\ .\ .\ .\ .}\}


\begin{center}
\textsc{Gauss} an \textsc{Schumacher}. Göttingen, 2. Januar 1836.
\end{center}

\mbox{.\ .\ .\ .\ .\ .\ .\ .} Die Pointe des Fehlers von Hrn. \textsc{Lübsen} besteht darin,
dass Euklids Geometrie falsch sein kann, ohne dass es in seiner
Construction bei
